\lecture{23}{Mon 18 Mar 2024 14:27}{C* algebra}

\subsection{C* Algebra and Continuous Functional Calculus}


\begin{definition}[C* Algebra]
	\label{defn:CAlgebra}
	A $C^*$-algebra is a Banach algebra endowed with a involution $*: A \to A$ which satisfies the following conditions:
	\begin{enumerate}
		\item (Involutive) $a^{* *} = a$
		\item (*-Linear) $(a + \lambda b)^* = a^* + \overline{\lambda} b^*$
		\item (sock-and-shoe) $(ab)^* = b^* a^*$
		\item (Normal?) $\|a^* a\| = \|a\|^2$
	\end{enumerate}
\end{definition}

\begin{definition}
	\label{defn:NormalUnitarySA}
	$a \in A$ is \textit{normal} if $a^* a = a a^*$.

	A normal element is \textit{self-adjoint} if $a^* = a$.

	A normal element is \textit{unitary} if $a a^* = a^* a =  1$.
\end{definition}
As we made explicit in the definition, self-adjoint and unitary elements are special cases of normal elements. As it will turn out later, those elements are precisely characterized by their spectrum (more precisely, spectral measure).

Now let's verify three basic facts about C* algebra. First, $*$ is continuous, and furthermore an isometry: $\|a\|^2 = \|a a^*\|\le \|a\| \|a^*\|$, so $\|a\| \le \|a^*\|$, by symmetry $\|a\| = \|a^*\|$.

Second is that involution is compatible with inverse: $(a^{-1})^* = \left( a^* \right) ^{-1}$. This is seen from $a^*(a^{-1})^* = (a^{-1} a)^* = 1^* = 1$. So, in particular, $*$ preserves invertability.

The third is that involution and inverse map the spectrum naturally: 
\begin{itemize}
	\item 
$\sigma_A(a^*) = \overline{\sigma_A(a)} $. This is simply seen: $(\lambda - a)$ invertible iff $(\lambda - a)^* = \overline{\lambda}  - a^*$ invertible. 

	\item $\sigma_A(a^{-1}) = \{\frac{1}{\lambda}: \lambda \in \sigma_A(a)\} $. This is seen by writing $(\lambda^{-1} - a^{-1}) = -\lambda^{-1} a^{-1} (\lambda - a)$.
\end{itemize}

This is an important special case for the more general \textit{spectral mapping theorem for continuous functional calculus.}




\begin{remark}
	For any element $a$, $a^* a$ is self-adjoint.
\end{remark}

First, let's see what normality says about an element's spectral radius.

\begin{proposition}
	\label{prop:NormalSpectralRadius}
	If $a \in A$ is normal, then $r_A(a) = \|a\|$.
\end{proposition}
\begin{proof}
	Assume first that $a$ is self-adjoint. Then $\|a^2\| = \|a^* a\| = \|a\|^2$. Iterating this gives $\|a^n\|^{1 / n} \to r_A(a) = \|a\|$.

	Now assume $a$ is normal. Then $a^* a$ is self-adjoint. By the previous argument, for $n = 2^r$,
	\[
		\|a^n\|^2 = \|a^n (a^n)^*\| = \|(a^* a)^n\| = \|a^* a\|^n = \|a\|^{2 n}
	.\] 
	Hence $\|a^n\| = \|a\|^n$. This gives $r_A(a) = \|a\|$.
\end{proof}

\begin{proposition}
	\label{prop:UnitSASpec}
	Let $A$ be a unital $C^*$-algebra.
	
	1. If $u$ is unitary then $\sigma_A(u) \subset \mathbb{T}$.
	
	2. If $u$ is self-adjoint, $\sigma_A(a) \subset \mathbb{R}$.
\end{proposition}
\begin{proof}
	(1) If $u$ is unitary, $u^* = u^{-1}$ is unitary as well. Now $\|u\|^2 = \|u^* u\| = 1 = r_A(u)$, and we also have $r_A(u^{-1}) = 1$ similarly. Now, for any $\lambda \in \sigma_A(a)$ we would have $|\lambda| \le 1$ and $|\frac{1}{\lambda}| \le 1$, so $|\lambda| = 1$.

	(2) If $a$ is self-adjoint, we define $u = e^{i a}$ using holomorphic functional calculus, and verify that $u$ is unitary. Then by the spectral mapping theorem for holomorphic functional calculus, for any $\lambda \in \sigma_A(a)$, we have $e^{i \lambda} \in \sigma_A(u) \subset \mathbb{T}$, so that $|e^{i \lambda}| = 1$, i.e. $\lambda $ is real.

	Take entire functions $f(z) = e^{i z}, g(z) = e^{- i z}$. By holomorphic functional calculus, $f(a) g(a) = g(a) f(a) = 1$. Thus $e^{- i a} = g(a) = f(a)^{-1} = (e^{i a})^{-1}$. We have
	\[
		(e^{i a})^* = \left( \sum_{k = 0}^\infty \frac{(i a)^k}{k!} \right) ^* = \sum \frac{(- i a)^k}{k!} = e^{- i a} = (e^{i a})^{-1}
	.\] 
	Note that we have used continuity of $*$. Thus, $u = e^{ia}$ is unitary.
\end{proof}

\begin{remark}
\textit{Positive Operators} are those elements $a \in A$ such that there exists $u \in A$ where $a = u^* u =: |u|^2$.

It can be shown that $a$ is positive if and only if $\sigma(a) \subset [0,\infty )$, and that positive elements have uniquely defined positive square roots.

We will omit the proofs for now, since square roots will be available to us once we have the powerful tool of spectral integration.
\end{remark}

Next, let's show that in unital $C^*$ algebra, the spectrum of an element does not depend on the ambient space.

\begin{lemma}
	\label{lem:InverseContinuous}
	Let $A$ be a unital Banach algebra. Then $a \mapsto a^{-1}$ defined on $\text{GL}(A)$ is a continuous operator.
\end{lemma}
\begin{proof}
	Let $a_n \to a$ in norm.
	Write
	\[
		\|a_n^{-1} - a^{-1}\|
		= \|a_n^{-1} a^{-1} \left( a - a_n \right) \|
		\le \|a_n^{-1}\| \|a^{-1}\| \|a - a_n\|
	.\] 
	Thus, it suffices to show $a_n^{-1}$ is norm-bounded as $a_n \to a$. Fix $r \in (0,1)$. For $\|a_n - a\| < \frac{r}{\|a^{-1}\|}$,
	\begin{align*}
		\|a_n^{-1}\| &= \|\sum_{k=0}^{\infty} (a^{-1} (a - a_n))^k a^{-1}\| \\
		&\leq \sum \|a^{-1} (a - a_n)\|^k \|a^{-1}\| \\
		&< \sum r^k \|a^{-1}\| \\
		&= \frac{\|a^{-1}\|}{1 - r} 
	.\end{align*}
\end{proof}


\begin{proposition}
	\label{prop:SpectrumIndependent}
	Let $B$ be a unital $C^*$ algebra, $A \subset B$ a norm-closed sub-algebra that contains $1_B$. Then $\sigma_A(a) = \sigma_B(a)$.
\end{proposition}
\begin{proof}
	If $a$ is invertible in $A$ then it is trivially invertible in $B$, so that $\sigma_B(a) \subset \sigma_A(a)$.

	Now we want to prove $\sigma_A(a) \subset \sigma_B(a)$. We have $a \in \text{GL}(A) \iff a^* a \in \text{GL}(A)$. Thus, it suffices to prove the statement for self-adjoint $a$. But then $\sigma_A(a), \sigma_B(a) \subset \mathbb{R}$.
	Take $\lambda \in \mathbb{R} \setminus  \sigma_B(a)$, so that $(\lambda - a)^{-1} \in B$. Now for each $n$, define $\lambda_n = \lambda + i \frac{1}{n}$. Since  $\lambda_n \not\in \sigma_A(a)$, $(\lambda_n - a)^{-1} \in A$. But $(\lambda_n - a)^{-1} \to (\lambda - a)^{-1}$ in norm, and $A$ is norm-closed, so $(\lambda - a)^{-1} \in A$. That is, $\lambda \not\in \sigma_A(a)$.
\end{proof}

\begin{proof}
	We give another direct proof which may provide more intuition. We reduce to the case of $a^* a$ and want to prove $\sigma_A(a^* a) \subset \sigma_B(a^* a)$. 

	By Proposition~\ref{prop:CstarCharacter}, $\omega(a^* a) = \overline{\omega(a)}  \omega(a) = |\omega(a)|^2 \ge 0$.

	By Proposition~\ref{prop:NormalSpectralRadius}, 
	\[
		\sigma_B(a^* a) \subset \sigma_A(a^* a) \subset [0,\|a^*a\|]
	.\] 
	Assume $a^* a \in \text{GL}(B)$. Then $0 \not\in \sigma_B(a^* a)$, thus
	\[
		\sigma_B(a^* a) \subset [\delta, \|a^* a\|]
	.\] 
	By holomorphic spectral mapping theorem,
	\[
		\sigma_B(\|a^* a\| - a^* a) = \|a^* a\| - \sigma_B(a^* a) \subset [0, \|a^* a\| - \delta]
	.\] 
	Now, $\|a^* a\| - a^* a$ is clearly self-adjoint, so
	\[
		\| \|a^* a\| - a^* a\| = r_B(\|a^* a\| - a^* a) \le  \|a^* a\| - \delta < \|a^* a\|
	.\] 

	Consider any $\omega \in \Omega(A)$. Since $\|\omega\| = 1$ (Proposition~\ref{prop:CharacterTopology}), 
	\[
		\omega\left( \|a^* a\| - a^* a \right) < \|a^* a\|
	.\] 
	But this implies $\omega(a^* a) > 0$, so $0 \not\in \sigma_A(a^* a)$, i.e. $a^* a \in \text{GL}(A)$.
\end{proof}
The main idea of this proof is that, although spectral information is not yet compatible between $A$ and $B$, they however share the same norm, so the inequality can be proved in $B$ using spectral information there, and then passed to $A$.

The theory for character spaces (Definition~\ref{defn:CharacterSpace}) and Galfand transform (Definition~\ref{defn:Gelftransform}) automatically holds for $C^*$ algebras, and further enjoy better properties:

\begin{proposition}
	\label{prop:CstarCharacter}
	Let $A$ be a unital commutative $C^*$-algebra. Then each element of $\Omega(A)$ is also a unital $*$-homomorphism.
\end{proposition}
\begin{proof}
	The unital homomorphism part is inherited. We prove that any $\omega \in \Omega(A)$ respects involution.

	Let $a \in $, write $a = x + i y$, where
	\[
		x = \frac{a + a^*}{2}, \quad y = \frac{a - a^*}{2 i}, \quad x, y \in A_{sa}
	.\] 
	Then for any $\omega \in \Omega(A)$ we have $\omega(x) \in \sigma_A(x) \subset \mathbb{R}$ and also $\omega(y) \in \mathbb{R}$, hence
	\[
		\omega(a^*) = \omega((x + i y)^*) = \omega(x - i y) = \omega(x) - i \omega(y) = \overline{\omega(x) + i \omega(y)}  = \overline{\omega(a)} 
	.\] 
\end{proof}

Now we have the following result that justifies focusing on $C^*$-algebras:

\begin{theorem}
	\label{thm:GelfandIsomorphism}
	Let $A $ be a unital and commutative $C^*$-algebra. Then the Gelfand transform $\Gamma: A \to C(\Omega(A))$ is an isometric $*$-isomorphism of $C^*$-algebras.
\end{theorem}
\begin{proof}
	Theorem~\ref{thm:GelfandProp} tells us $\|\Gamma(a)\| = r_A(a)$. Since $A$ is commutative, $a$ is normal, hence $r_A(a) = \|a\|$ (Proposition~\ref{prop:NormalSpectralRadius}), hence $\|\Gamma(a)\| = \|a\|$.

	Proposition~\ref{prop:CstarCharacter} tells us $\Gamma(a^*) = \overline{\Gamma(a)} $. Therefore, $\Gamma$ is an isometric $*$-homomorphism.

	Finally we need to show $\Gamma$ is an isomorphism, but since it is isometric, we only need to show $\Gamma$ is surjective. To do so, we apply the Stone-Weierstrass Theorem to $\Gamma(A) \subset C(\Omega(A))$.
	\begin{enumerate}
		\item $\Gamma(1) = 1 \in C(\Omega(A))$.
		\item Pick $\omega_1 \neq \omega_2 \in \Omega(A)$. Then $\omega_1(a) \neq \omega_2(a)$ for $a \in A$. But then $\Gamma(a)$ separates $\omega_1$ and $\omega_2$.
		\item For every $\Gamma(a) \in \Gamma(A)$, $\overline{\Gamma(a)}  = \Gamma(a^*) \in \Gamma(A)$.
	\end{enumerate}
	Hence, $\Gamma(A)$ is dense in $C(\Omega(A))$. Since $A$ is complete and $\Gamma$ is isometric, $\Gamma(A)$ is complete, thus $\Gamma(A) = C(\Omega(A))$.
\end{proof}

If $A$ is a unital $C^*$-algebra and $a \in A$ is a normal element, we can consider the $C^*$-algebra generated by $1$ and $a$. It has the following form:
\[
	C^* \{1, a\}  = \overline{span}\left( a^m a^{* n}, m, n \ge 0 \right)  
.\] 
Note that we are taking norm closure of linear span. Moreover, $C^* \{1, a\} $ is a unital commutative $C^*$-algebra.

In view of Proposition~\ref{prop:SpectrumIndependent}, we may consider the subalgebra $C^* \{1, a\} $ for any normal element $a$ in some unital $C^*$-algebra. This approach turns out to be extremely helpful:

\begin{theorem}
	\label{thm:CharacterSpectrumHomeo}
	Let $A = C^* \{1, a\} $ where $a$ is a normal element. Then
	\[
		\Omega(A) \ni \omega \mapsto \omega(a) \in \sigma(a)
	\] 
	is a homeomorphism.
\end{theorem}
\begin{proof}
	Denote the map by $h$.
	If $\omega_\alpha \to \omega$ in weak-* topology, then for each $a \in A$, $\omega_\alpha(a) \to \omega(a)$ in norm. Thus $h$ is continuous.

	Since $\sigma(a) = \{\omega(a): \omega \in \Omega(A)\} $ (Proposition~\ref{prop:CharacterSpectrum}), $h$ is surjective.

	Now we verify $h$ is injective. Take $\omega_1, \omega_2 \in \Omega(A)$, $\omega_1(a) = \omega_2(a)$. Then $\omega_1(a^*) = \omega_2(a^*)$. Thus, $\omega_1, \omega_2$ agree on $*$-algebra$\{1, a\} $. By continuity, they agree on $C^*\{1, a\}  = A$.
\end{proof}

Using this homeomorphism $h: \Omega(A) \to \sigma(a)$, we can induce a $*$-isomorphism $h^*: C(\sigma(a)) \to C(\Omega(A))$. Combine this with the Gelfand transform, we obtain the 

\begin{theorem}[Continuous Functional Calculus]
	\label{thm:ContinuousFunctionalCalculus}
	Let $B$ be a unital $C^*$-algebra and let $a \in B$ be a normal element. Then there is a unital $*$-isomorphism
	\[
		\varphi: C(\sigma(a)) \to C^*\{1,a\} = A 
	\] 
	such that $\varphi(z) = a$.
\end{theorem}
\begin{proof}
	The desired $*$-isomorphism is built directly by
	\[
		C(\sigma(a)) \xrightarrow{h^*} C(\Omega(A))
		\xleftarrow{\Gamma} C^*\{1, a\}  = A
	.\] 
	By Theorem~\ref{thm:GelfandIsomorphism} and Theorem~\ref{thm:CharacterSpectrumHomeo}, we indeed have an isometric $*$-isomorphism.

	Now, we want to verify $\varphi(z) = a$. We have $\Gamma(\varphi(z)) = h^*(z)$, $z(\lambda) = \lambda$. Evaluate both sides on $\omega \in \Omega(A)$ :
	\[
		h^*(z)(\omega) = (z \circ h)(\omega) = h(\omega) = \omega(a)
	.\] 
	On the other hand,
	\[
		\Gamma(\varphi(z))(\omega) = \omega(\varphi(z))
	.\] 
	Therefore $\omega(a) = \omega(\varphi(z))$. This is true for all $\omega \in \Omega(A)$, so $\varphi(z) = a$.
\end{proof}
\ask{Justify the last step more carefully. Why $\Omega(A)$ separates? Gelfand-Naimark theorem?}


Let's discuss some properties of the continuous functional calculus.

\begin{corollary}(Functional Mapping Theorem)

	\label{cor:ContSpectralMapping}
	$\sigma(f(a)) = f(\sigma(a))$.
\end{corollary}

The continuous functional calculus is functorial:
\begin{proposition}
	\label{prop:ContFunctorial}
	 Let  $\pi:A \to B$ be a unital $*$-homomorphism, let $a \in A$ be a normal element. Then
	\begin{enumerate}
		\item $f(\pi(a)) = \pi(f(a))$.
		\item If $f \in C(\sigma(a)), g \in C(f(\sigma(a)))$, then $g(f(a)) = (g \circ f) (a)$.
	\end{enumerate}
\end{proposition}
Moreover, $\varphi$ is an isometry, as implied by the next proposition:
\begin{proposition}
	\label{prop:StarHomoIsometry}
	Let $\pi: A \to B$ be a unital $*$-homomorphism. Then
	\begin{enumerate}
		\item $\|\pi(a)\| \le  \|a\|$ for all $a \in A$,
		\item $\pi$ injective iff $\|\pi(a)\| = \|a\|$ for all $a \in A$.
	\end{enumerate}
\end{proposition}



The next topic of study is \textbf{diagonalization.}











